%---------------------------------------------------------------------------------------------------------------------------- 
\graphicspath{{Parti/P03-Meccanica_SR/C03-Dualita_CR/Figure/}}
\chapter{Dualismo tra cinematica e statica del corpo rigido}\label{cap:dualismo_CR}
%---------------------------------------------------------------------------------------------------------------------------- 


\begin{sommario}
	Si mette in luce il dualismo tra il campo degli spostamenti	rigidi infinitesimi e il campo dei momenti: le due leggi di	distribuzione hanno la stessa struttura affine, e a ogni concetto cinematico corrisponde un concetto statico duale.	Si specializza il confronto al caso piano, evidenziando lo	scambio di ruolo tra scalare e vettore che trasforma il centro di rotazione nell'asse di sollecitazione. Si introduce il lavoro virtuale come forma bilineare che
	accoppia le grandezze statiche e cinematiche, e se ne	calcola l'espressione per forze concentrate, coppie e
	forze distribuite; per queste ultime si mostra	l'equivalenza tra la scrittura in termini di grandezze
	generalizzate e quella in termini di risultante applicata	nel baricentro delle forze. Si conclude mostrando come la	condizione di lavoro virtuale nullo per ogni spostamento	rigido sia equivalente alle equazioni di equilibrio del	corpo rigido.
\end{sommario}

\section{Dualismo cinematico-statico}\label{sec:dualismo}\index{dualita@dualità!cinematica--statica|textbf}\index{dualismo|see{dualità}}

Le due leggi di distribuzione --- quella degli spostamenti
rigidi infinitesimi e quella dei momenti --- hanno la stessa
struttura:
\begin{equation}\label{eq:dualismo}
	\begin{aligned}
		\text{campo cinematico:} &\qquad
		\ub(P) = \ub(O)
		+ \bm{\upvartheta} \times \overrightarrow{OP}, \\[4pt]
		\text{campo statico:} &\qquad
		\bm{\upmu}(P) = \bm{\upmu}(O)
		+ \fb \times \overrightarrow{OP}.
	\end{aligned}
\end{equation}
La corrispondenza si estende a tutte le proprietà derivate
nei capitoli precedenti:

\begin{center}
	\renewcommand{\arraystretch}{1.4}
	\begin{tabular}{ccc}
		\toprule
		\textbf{Campo cinematico} & $\longleftrightarrow$
		& \textbf{Campo statico} \\
		\midrule
		spostamento $\ub(P)$
		& & momento $\bm{\upmu}(P)$ \\
		spostamento di $O$: $\ub(O)$
		& & momento al polo: $\bm{\upmu}(O)$ \\
		rotazione $\bm{\upvartheta}$
		& & risultante $\fb$ \\
		invariante $\bm{\upvartheta}\cdot\ub(O)$
		& & invariante $\fb\cdot\bm{\upmu}(O)$ \\
		asse centrale degli spostamenti
		& & asse centrale delle forze \\
		distanza $|\ub_\perp(O)|/|\bm{\upvartheta}|$
		& & distanza $|\bm{\upmu}_\perp(O)|/|\fb|$ \\
		\bottomrule
	\end{tabular}
\end{center}

\subsection{Il caso piano}

Nel piano la corrispondenza si specializza:

\begin{center}
	\renewcommand{\arraystretch}{1.4}
	\begin{tabular}{ccc}
		\toprule
		\textbf{Cinematica} & $\longleftrightarrow$
		& \textbf{Statica} \\
		\midrule
		spostamento di $O$: $\ub(O)$
		& & momento al polo: $\mu_z(O)$ \\
		rotazione $\vartheta$
		& & risultante $\fb$ \\
		centro di rotazione (punto)
		& & asse di sollecitazione (retta) \\
		$|\overrightarrow{OC}| = |\ub(O)|/|\vartheta|$
		& & $|\overrightarrow{OC}_\perp|
		= |\mu_z(O)|/|\fb|$ \\
		$\overrightarrow{OC}
		= (\ab_z \times \ub(O))/\vartheta$
		& & $\overrightarrow{OC}_\perp
		= -\mu_z(O)\,(\ab_z \times \fb)/|\fb|^2$ \\
		traslazione pura: $C$ all'infinito
		& & coppia pura: asse all'infinito \\
		baricentro dei centri di rotazione
		& & baricentro dei punti di applicazione \\
		\bottomrule
	\end{tabular}
\end{center}

\begin{oss}
	Il dualismo è perfetto nella struttura algebrica ma
	presenta una differenza geometrica nel piano: la rotazione
	$\vartheta$ è uno scalare e $\ub(O)$ è un vettore a due
	componenti, quindi $\ub(C)=\zerovb$ impone due equazioni
	in due incognite e determina un \textit{punto}. La
	risultante $\fb$ è un vettore a due componenti e
	$\mu_z(O)$ è uno scalare, quindi $\mu_z(C)=0$ impone una
	equazione in due incognite e determina una \textit{retta}.
	Lo scambio di ruolo fra scalare e vettore produce lo
	scambio fra punto e retta.
\end{oss}



\begin{oss}
	I due campi differiscono per natura: la formula
	cinematica $\ub(P) = \ub(O) + \bm{\upvartheta}\times\overrightarrow{OP}$ \textit{è} il campo --- esprime direttamente
	lo spostamento di ogni punto. La formula statica $\bm{\upmu}(P) = \bm{\upmu}(O)
	+ \fb\times\overrightarrow{OP}$ è un
	\textit{teorema} ricavato dalla definizione di momento
	risultante, in cui il polo $O$ ha un ruolo privilegiato.
	Il dualismo risiede nel fatto che entrambi i campi
	obbediscono alla stessa legge affine.
\end{oss}


\section{Lavoro virtuale}\label{sec:lavoro_virtuale}

Il dualismo trova la sua espressione più concreta nel
\textit{lavoro virtuale}\index{lavoro virtuale|textbf}\nomenclature[L]{$\Wc$}{lavoro virtuale}. Dato un sistema di forze applicate
a un corpo rigido e uno spostamento rigido infinitesimo,
il lavoro virtuale è la somma dei lavori elementari delle
singole forze nei rispettivi spostamenti.

\subsection{Forze concentrate}

Consideriamo $n$ forze $\fb_k$ applicate nei punti $P_k$.
Lo spostamento di $P_k$ è
$\ub_k = \ub(O) + \bm{\upvartheta} \times
\overrightarrow{OP_k}$, dunque
\begin{equation}
	\Wc = \sum_{k=1}^{n} \fb_k \cdot \ub_k
	= \sum_{k=1}^{n} \fb_k \cdot \ub(O)
	+ \sum_{k=1}^{n} \fb_k \cdot
	(\bm{\upvartheta} \times \overrightarrow{OP_k}).
\end{equation}
Sfruttando la proprietà del prodotto misto
$\ab \cdot (\bb \times \cb) = (\cb \times \ab) \cdot \bb$:
\begin{equation}
	\Wc = \Bigl(\sum_k \fb_k\Bigr) \cdot \ub(O)
	+ \Bigl(\sum_k \overrightarrow{OP_k}
	\times \fb_k\Bigr) \cdot \bm{\upvartheta}.
\end{equation}
Si riconoscono risultante e momento risultante:
\begin{equation}\label{eq:lavoro_virtuale}
	\boxed{\Wc = \fb \cdot \ub(O)
		+ \bm{\upmu}(O) \cdot \bm{\upvartheta}.}
\end{equation}
Il lavoro dipende unicamente dalle grandezze globali del
sistema di forze e dallo spostamento rigido: risultante,
momento risultante, traslazione del polo e rotazione.

\subsection{Coppie concentrate}

Una coppia $\bm{\upmu}_j$ (vettore assiale libero) applicata
al corpo rigido compie un lavoro virtuale
\begin{equation}
	\Wc_j = \bm{\upmu}_j \cdot \bm{\upvartheta}.
\end{equation}
Per verificarlo, si scrive la coppia come due forze
$\fb_1 = \fb$ e $\fb_2 = -\fb$ applicate in $P_1$ e $P_2$,
con $\bm{\upmu}_j = \overrightarrow{P_2P_1} \times \fb$.
Il lavoro delle due forze è
\begin{equation}
	\Wc_j = \fb \cdot \ub(P_1) - \fb \cdot \ub(P_2)
	= \fb \cdot [\ub(P_1) - \ub(P_2)]
	= \fb \cdot (\bm{\upvartheta} \times
	\overrightarrow{P_2P_1})
	= \bm{\upmu}_j \cdot \bm{\upvartheta}.
\end{equation}
Il lavoro di una coppia dipende solo dalla rotazione, non
dalla traslazione del polo: una coppia non compie lavoro
in una traslazione pura.

Per un sistema misto di forze concentrate e coppie,
la~\eqref{eq:lavoro_virtuale} si generalizza a
\begin{equation}\label{eq:lavoro_misto}
	\Wc = \fb \cdot \ub(O)
	+ \bm{\upmu}(O) \cdot \bm{\upvartheta},
\end{equation}
dove $\bm{\upmu}(O)$ include ora anche il contributo
diretto delle coppie secondo la~\eqref{eq:riduzione_generalizzata}.


\begin{oss}
	L'espressione~\eqref{eq:lavoro_virtuale} può essere
	scritta in forma compatta come
	\begin{equation}
		\Wc = \bm{f}^\top \bm{u},
	\end{equation}
	dove
	$\bm{f} = \{X\; Y\; Z\;
	\mu_x(O)\; \mu_y(O)\; \mu_z(O)\}^\top$
	è il vettore delle grandezze statiche e
	$\bm{u} = \{u(O)\; v(O)\; w(O)\;
	\vartheta_x\; \vartheta_y\; \vartheta_z\}^\top$
	è il vettore delle grandezze cinematiche. Le due
	famiglie di grandezze sono \textit{duali} nel senso del
	\CAPref{app:applicazionilin_dualita}: il loro prodotto
	scalare ha il significato fisico di lavoro.
\end{oss}

\subsection{Forze distribuite}

Per un carico distribuito $\fb_L(P)$ lungo una linea, il
lavoro virtuale è
\begin{equation}
	\Wc = \int_L \fb_L(P) \cdot \ub(P)\,\dd L.
\end{equation}
Sostituendo la formula fondamentale
$\ub(P) = \ub(O) + \bm{\upvartheta} \times
\overrightarrow{OP}$ e procedendo come nel caso concentrato:
\begin{equation}
	\Wc = \fb \cdot \ub(O)
	+ \bm{\upmu}(O) \cdot \bm{\upvartheta},
\end{equation}
con
\begin{equation}
	\fb = \int_L \fb_L\,\dd L, \qquad
	\bm{\upmu}(O) = \int_L \overrightarrow{OP}
	\times \fb_L\,\dd L.
\end{equation}
La forma è identica al caso concentrato: la distinzione
fra forze concentrate e distribuite scompare a livello
di grandezze globali.


\paragraph{Scrittura alternativa.}
Nel caso piano, quando $\fb \neq \zerovb$, il sistema è
sempre equivalente a una singola forza $\fb$ applicata
sull'asse di sollecitazione: l'invariante
$\fb \cdot \bm{\upmu}(O)$ è identicamente nullo perché
$\fb$ giace nel piano e $\bm{\upmu}(O)$ è ortogonale ad
esso. Scegliendo come polo un punto $C$ sull'asse di
sollecitazione, dove $\bm{\upmu}(C) = \zerovb$,
la~\eqref{eq:lavoro_virtuale} diventa
\begin{equation}\label{eq:lavoro_risultante}
	\Wc = \fb \cdot \ub(C) +
	\underbrace{\bm{\upmu}(C)}_{=\,\zerovb}
	\cdot \bm{\upvartheta}
	= \fb \cdot \ub(C).
\end{equation}
Il lavoro virtuale del sistema coincide con quello della
sola forza risultante applicata in $C$. Questo risultato
vale per un corpo rigido di forma qualsiasi.

Nel caso particolare di un carico distribuito
$\fb_L(x) = -p(x)\,\ab_y$ con $p(x) \geq 0$, la risultante
è $\fb = -|\fb|\,\ab_y$ con $|\fb| = \int_0^\ell p(x)\,\dd x$,
applicata nel baricentro $C$ del diagramma di carico. Il
lavoro virtuale diventa
\begin{equation}
	\Wc = \fb \cdot \ub(C) = -|\fb|\,v(C),
\end{equation}
dove $v(C)$ è la componente dello spostamento di $C$ nella
direzione del carico: il lavoro è il prodotto fra l'area
del diagramma di carico e lo spostamento nella direzione
del carico nel suo baricentro. Un analogo risultato vale
per un carico $\fb_L(y) = -p(y)\,\ab_x$, con $u(C)$
al posto di $v(C)$.



\section{Equilibrio e lavoro virtuale}\index{equilibrio}
\label{sec:equilibrio_LV}

La condizione di equilibrio del corpo rigido --- già
richiamata nel \CAPref{cap:statica_CR} come dato dei corsi
propedeutici --- trova nel lavoro virtuale una formulazione
alternativa e equivalente.

\begin{teorema}[Principio dei lavori virtuali]\index{principio dei lavori virtuali|textbf}\index{PLV|see{principio dei lavori virtuali}}
	Un corpo rigido è in equilibrio se e solo se il lavoro
	virtuale delle forze esterne è nullo per ogni
	spostamento rigido infinitesimo:
	\begin{equation}\label{eq:PLV}
		\Wc = \fb \cdot \ub(O)
		+ \bm{\upmu}(O) \cdot \bm{\upvartheta} = 0
		\qquad \forall\,\ub(O),\,\bm{\upvartheta}.
	\end{equation}
\end{teorema}
La dimostrazione è diretta. Poiché $\ub(O)$ e
$\bm{\upvartheta}$ sono indipendenti --- rappresentano
rispettivamente la traslazione e la rotazione, che
costituiscono i sei gradi di libertà indipendenti del
corpo rigido --- è lecito sceglierli separatamente.
Scegliendo $\bm{\upvartheta} = \zerovb$ e $\ub(O)$
arbitrario, la~\eqref{eq:PLV} richiede
$\fb \cdot \ub(O) = 0$ per ogni $\ub(O)$, dunque
$\fb = \zerovb$. Scegliendo poi $\ub(O) = \zerovb$ e
$\bm{\upvartheta}$ arbitrario, la~\eqref{eq:PLV} richiede
$\bm{\upmu}(O) \cdot \bm{\upvartheta} = 0$ per ogni
$\bm{\upvartheta}$, dunque $\bm{\upmu}(O) = \zerovb$.
Le equazioni di equilibrio~\eqref{eq:equilibrio_CR}
sono dunque equivalenti alla~\eqref{eq:PLV}.


\begin{oss}
	Il principio dei lavori virtuali non è una nuova legge
	fisica: è una riformulazione delle equazioni di
	equilibrio in termini di lavoro. Il suo valore sta
	nella generalità: si applica senza modifiche a sistemi
	di corpi connessi, a strutture con vincoli, e --- nella
	sua estensione ai corpi deformabili --- diventa il
	fondamento della meccanica computazionale. Il filo
	conduttore che lega la dualità algebrica introdotta nel
	\CAPref{app:applicazionilin_dualita} al lavoro virtuale
	sarà sviluppato sistematicamente nel seguito del testo.
\end{oss}


\begin{oss}
	Nel caso piano la~\eqref{eq:lavoro_virtuale} si specializza in
	\begin{equation}\label{eq:lavoro_piano}
		\Wc = X\,u(O) + Y\,v(O) + \mu(O)\,\vartheta,
	\end{equation}
	e la~\eqref{eq:PLV} diventa
	\begin{equation}
		X\,u(O) + Y\,v(O) + \mu(O)\,\vartheta = 0
		\qquad \forall\,u(O),\,v(O),\,\vartheta,
	\end{equation}
	da cui le tre equazioni di equilibrio scalari $X = 0$, $Y = 0$, $\mu(O) = 0$. I tre parametri cinematici $u(O)$, $v(O)$, $\vartheta$ e i tre parametri statici $X$, $Y$, $\mu(O)$ sono duali nel senso della~\eqref{eq:lavoro_piano}: il loro prodotto scalare è il lavoro virtuale.
\end{oss}


%-------------------------------------------------------------------------------
\section{Esercizi proposti}
%-------------------------------------------------------------------------------

\begin{esercizio}[Dualismo cinematica-statica nel piano]
	\label{ex:C03_dualismo_piano}
	
	Un corpo rigido piano è soggetto, contemporaneamente, a uno spostamento rigido infinitesimo e a un sistema di forze. Lo spostamento del polo è $\ub(O) = (0.4\;\,-0.2)^\top$ con rotazione $\vartheta = 0.05$\,rad. La risultante delle forze è $\fb = (3\;\,-1)^\top$\,kN e il momento rispetto a $O$ è $\mu(O) = 2$\,kN$\cdot$m.
	(a) Determinare il centro di rotazione $C_R$ e l'equazione dell'asse di sollecitazione.
	(b) Compilare la tabella di corrispondenza cinematica-statica con i valori numerici.
	(c) Calcolare il lavoro virtuale del sistema di forze sullo spostamento, usando la~\eqref{eq:lavoro_piano}.
	
	\risultato (a) $C_R$: $x(C_R) = -v(O)/\vartheta = 0.2/0.05 = 4$, $y(C_R) = u(O)/\vartheta = 0.4/0.05 = 8$, dunque $C_R = (4,\,8)$. Asse di sollecitazione: $\mu(O) + X\,y - Y\,x = 0$, ovvero $2 + 3y + x = 0$, equivalente a $x + 3y + 2 = 0$. (b) Tabella di corrispondenza:
	\begin{center}
		\begin{tabular}{ll}
			\toprule
			\textbf{Cinematica} & \textbf{Statica} \\
			\midrule
			$\ub(O) = (0.4\;\,-0.2)^\top$ & $\mu(O) = 2$\,kN$\cdot$m \\
			$\vartheta = 0.05$\,rad & $\fb = (3\;\,-1)^\top$\,kN \\
			$C_R = (4,\,8)$ & asse di sollecitazione: $x + 3y + 2 = 0$ \\
			\bottomrule
		\end{tabular}
	\end{center}
	Si osservi lo \textit{scambio di ruolo scalare-vettore}: nel piano $\bm{\upvartheta}$ (cinematica) è scalare e $\fb$ (statica) è vettore; viceversa $\ub(O)$ (cinematica) è vettore e $\mu(O)$ (statica) è scalare. (c) $\Wc = X\,u(O) + Y\,v(O) + \mu(O)\,\vartheta = 3 \cdot 0.4 + (-1)(-0.2) + 2 \cdot 0.05 = 1.2 + 0.2 + 0.1 = 1.5$\,kJ.
	
\end{esercizio}

\begin{esercizio}[Dal lavoro virtuale all'equilibrio]
	\label{ex:C03_LV_equilibrio}
	
	Un corpo rigido piano è soggetto a un sistema di forze con risultante $\fb = (X\;\,Y)^\top$ e momento $\mu(O)$ rispetto a un polo $O$. Si supponga che il lavoro virtuale del sistema, $\Wc = X\,u(O) + Y\,v(O) + \mu(O)\,\vartheta$, sia nullo per \textit{ogni} spostamento rigido virtuale $\{u(O),\,v(O),\,\vartheta\}$.
	(a) Mostrare che ciò implica le tre equazioni di equilibrio scalari $X = 0$, $Y = 0$, $\mu(O) = 0$.
	(b) Commentare il legame con il principio dei lavori virtuali~\eqref{eq:PLV}.
	
	\risultato (a) Poiché i tre parametri cinematici $u(O)$, $v(O)$, $\vartheta$ sono indipendenti --- corrispondono ai tre gradi di libertà del corpo rigido nel piano --- è lecito sceglierli separatamente. Ponendo $u(O) = 1$, $v(O) = \vartheta = 0$, la condizione $\Wc = 0$ dà $X = 0$. Ponendo $v(O) = 1$, $u(O) = \vartheta = 0$, dà $Y = 0$. Ponendo $\vartheta = 1$, $u(O) = v(O) = 0$, dà $\mu(O) = 0$. Le tre equazioni di equilibrio scalari sono dunque condizioni necessarie. La sufficienza è immediata: se $X = Y = \mu(O) = 0$, allora $\Wc = 0$ identicamente. (b) Questo è esattamente il \textit{principio dei lavori virtuali}~\eqref{eq:PLV}, specializzato al caso piano. Nello spazio le equazioni di equilibrio diventano sei (tre per la risultante $\fb$, tre per il momento $\bm{\upmu}(O)$), ed equivalgono alla condizione $\Wc = \fb \cdot \ub(O) + \bm{\upmu}(O) \cdot \bm{\upvartheta} = 0$ per ogni $\ub(O)$ e ogni $\bm{\upvartheta}$. Il PLV non è dunque una nuova legge fisica, ma una riformulazione bilineare delle equazioni di equilibrio.
	
\end{esercizio}



	
